\section{Discussion}\label{sec:discussion}

\noindent \textbf{Scope of detection.} \toolname{} is designed to detect malicious behaviors that a model artifact executes against the host system, such as data exfiltration. Its detection pipeline relies on observing the actual runtime interactions between the model process and the operating system. \toolname{} could not address attacks that target the model itself without affecting the host system. In particular, model-level threats such as neural backdoor injection fall outside our detection scope. These attacks embed malicious functionality into the model's learned parameters, and at inference time the compromised model performs only normal tensor computations through standard framework operations. So these attacks generate no anomalous system calls or suspicious API invocations that \toolname{} relies upon as its primary detection signal. 

Consequently, these attacks are invisible to \toolname{}'s cross-layer monitoring pipeline. Such attacks compromise the model's prediction integrity rather than its runtime behavior, and their malicious effects manifest as incorrect or biased outputs rather than as anomalous system calls or API invocations. Defending against such threats requires complementary techniques such as model provenance verification, weight integrity checking, and robust training, which are orthogonal to and compatible with \toolname{}'s host-level defense.


\section{Conclusion}\label{sec:conclusion}

Malicious model artifacts present a growing threat to the AI supply chain: attackers can embed harmful payloads within serialized models that execute during loading or inference, evading traditional static scanners by using legitimate framework APIs. In this paper, we propose \toolname{}, a hybrid runtime defense system that combines non-bypassable kernel-level system-call monitoring with LLM-assisted Python semantic instrumentation to detect malicious behaviors in model artifacts. Instead of relying solely on static signature matching, \toolname{} captures the actual runtime execution evidence through a two-stage architecture: an unsupervised Autoencoder prefilter that surfaces anomalous low-level operations, and a cross-layer LLM reasoning engine that aligns these operations with their high-level semantic context to determine malicious intent. We evaluate \toolname{} on the MalHug Dataset and advanced synthetic attacks. The results show that \toolname{} achieves near-perfect detection, while introducing an average runtime overhead of only 5.69\%, demonstrating that \toolname{} provides a practical, deployable defense for the AI model supply chain.